\begin{abstract}
Song libraries are searched by people who need to name a bird from a field recording. We ask whether retrieval scores on such a library measure matching sounds to species or only matching recordings to the archive they came from. We built a small benchmark of forty recordings and tested three audio models against a baseline that sees only the archive label. One model beat that baseline, and only on a minority of queries.
\end{abstract}
\section{Introduction}
Field guides now ship with audio search. A listener records a song and the tool returns candidate species. Scores on public benchmarks are read as evidence that such tools match sounds. Archives differ in microphones, background noise and editing, and a model can learn those differences instead of the birds.
\section{Results}
The archive baseline named 12 species correctly at the first rank. The effect held only for the forest archive.
